\documentclass[11pt]{article}
\usepackage[margin=1in]{geometry}
\usepackage{amsmath,amssymb,booktabs,array,longtable}
\usepackage[dvipsnames]{xcolor}
\usepackage{tcolorbox}
\usepackage{hyperref}
\hypersetup{colorlinks=true,linkcolor=RoyalBlue,urlcolor=RoyalBlue,citecolor=RoyalBlue}
\usepackage{enumitem}
\sloppy

\newtcolorbox{honestbox}{colback=Orange!8,colframe=Orange!75!black,boxrule=0.6pt,arc=2pt,
  title=\textbf{Honest scope},fonttitle=\bfseries,left=4pt,right=4pt,top=3pt,bottom=3pt}
\newtcolorbox{decisionbox}{colback=RoyalBlue!6,colframe=RoyalBlue!65!black,boxrule=0.6pt,arc=2pt,
  title=\textbf{Decision rule},fonttitle=\bfseries,left=4pt,right=4pt,top=3pt,bottom=3pt}

\title{\textbf{Holonet Demonstrator: the Parity Control Arm}\\[0.3em]
\large Turning the $W(3,3)$ Contextuality Test into a Two-Arm Discriminator}
\author{Wil Dahn \quad\textbar\quad \texttt{github.com/wilcompute/W33-Theory}}
\date{June 2026}

\begin{document}
\maketitle

\begin{abstract}
The Holonet Demonstrator Protocol v1 measures the contextual fraction of the $q=3$ Witting fabric and
predicts $\mathrm{CF}=1/10$.  This note adds the missing \emph{positive control}.  The parity law --- a
generalized quadrangle $W(q)$ has an ovoid, hence a noncontextual model, if and only if $q$ is even ---
predicts that the \emph{same class} of single-photon apparatus, re-pointed at an \emph{even}-order
analyzer map ($W(2)$ with $15$ rays or $W(4)$ with $85$ rays), must read $\mathrm{CF}=0$.  Running both
arms turns a one-sided test (``observe $1/10$, or kill'') into a \emph{discriminating} measurement:
contextuality must appear on the odd fabric and vanish on the even fabric, with the contrast fixed by
the geometry, not by any tunable parameter.  A systematic that could fake a violation would have to
respect the parity of the field order to survive both arms.  The control arm ships with an explicit
noncontextual model: the ovoid assignment constructed and verified in
\texttt{analysis/w33\_ovoid\_construct.py}.
\end{abstract}

\section{The two arms}
The substrate is the symplectic generalized quadrangle $W(q)=\mathrm{GQ}(q,q)$, whose points are the
measurement rays and whose totally-isotropic lines are the $n=(q+1)(q^2+1)$ measurement contexts.  The
parity law (witness \texttt{analysis/w33\_audit\_qscan.py}) gives the contextual fraction in closed
form,
\[
  \mathrm{CF}(q)=
  \begin{cases}
    0, & q\ \text{even},\\[2pt]
    \dfrac{1}{q^2+1}, & q\ \text{odd},
  \end{cases}
\]
so the experiment has two arms with opposite, geometry-forced predictions:

\begin{center}
\begin{tabular}{lcccc}
\toprule
Arm & Fabric & rays $n$ & contexts & predicted $\mathrm{CF}$\\
\midrule
positive (v1) & $W(3)$ & $40$ & $40$ & $1/10$\\
control       & $W(2)$ & $15$ & $15$ & $0$\\
control       & $W(4)$ & $85$ & $85$ & $0$\\
\bottomrule
\end{tabular}
\end{center}

On the even-order arm an \emph{ovoid} exists: a set of $q^2+1$ rays (five for $W(2)$, seventeen for
$W(4)$) meeting every context exactly once.  Assigning the value $1$ to those rays and $0$ to all others
is a single, context-independent truth assignment that satisfies all $n$ contexts simultaneously --- a
working noncontextual hidden-variable model.  Its existence is precisely why the even arm must measure
no contextuality.

\section{The explicit control model}
The control prediction is not an assertion; it is a constructed object.  The witness
\texttt{analysis/w33\_ovoid\_construct.py} builds the ovoid for $q=2$ and $q=4$ and verifies it three
independent ways: it has size $q^2+1$, it meets every context in exactly one ray, and its rays are
pairwise non-collinear (a cap).  The same witness confirms that the odd order $q=3$ admits \emph{no}
such assignment (the best leaves four of forty contexts unsatisfiable), which is the obstruction the
positive arm detects.  The explicit $W(2)$ ovoid is
\[
  \{0100,\ 1000,\ 1101,\ 1110,\ 1111\},
\]
five rays, one per context; the $W(4)$ ovoid has seventeen rays and is written to
\texttt{data/w33\_ovoid\_construct.json}.

\section{Two notions of contextuality (what the control arm does and does not test)}
Two exact results sharpen the scope of the control prediction, both computed in the repository.

First, \emph{the even fabric is not globally classical}.  Labelling the fifteen points of $W(2)$ by the
fifteen nontrivial two-qubit Pauli operators, every line is a commuting triple whose product is $\pm I$,
and the noncontextual $\pm1$ assignment system is \emph{unsatisfiable}: the witness
\texttt{analysis/w33\_doily\_mermin.py} computes the fifteen signs as $4\times4$ matrix products, proves
unsatisfiability by exact $F_2$ elimination, and extracts the minimal certificate --- six lines, every
point covered an even number of times, an odd number of $-I$ lines: a Mermin--Peres magic square inside
the doily.  So $W(2)$ is \emph{sign}-contextual.  What the parity law kills on even $q$ is specifically
the \emph{selection} contextuality --- the exactly-one-click-per-context excess that the
contextual-fraction estimators measure --- for which the ovoid is a perfect noncontextual model.  The
$q=3$ fabric is contextual in \emph{both} senses.  The two-arm contrast is therefore a statement about
one declared observable, the CF statistic, and both arms must measure exactly that statistic.

Second, \emph{the control arm provably needs a different measurement class}.  A complete-basis
realization of $W(q)$ maps each context to an orthonormal basis, hence lives in $\mathbb C^{q+1}$.  For
$q=2$ this is impossible: a non-collinear pair of rays spans two dimensions of $\mathbb C^3$, leaving a
one-dimensional orthocomplement that contains exactly one ray, while the geometry demands $\mu=3$
distinct common neighbours there (\texttt{analysis/w33\_realization\_dimension.py} verifies the counts).
So the doily cannot be interrogated by the $q=3$ arm's complete von Neumann measurements in its matching
dimension; its exactly-one-click statistic must be realized by a different measurement class, as the
honest-scope box below already concedes.  The same counting at $q=3$ poses no obstruction --- and the
Witting realization exists --- so $q=3$ is the \emph{smallest} order the one-photon $\mathbb C^4$
apparatus can realize at all, and by the parity law the smallest contextual one: two independent
forcings of the same order.

\section{Statistics and decision rule}
Each arm uses the same raw-shot schema, validator, and estimator as v1.  Distinguishing a contextual
fraction $p$ from $0$ at $k\sigma$ needs, in the binomial planning approximation,
\[
  N \ge \frac{k^2\,p(1-p)}{p^2},
\]
so the positive arm targets $N\approx 225$ valid coincidences for $5\sigma$ at $p=1/10$.  The control
arm instead \emph{bounds $\mathrm{CF}$ from above}: it must show no statistically significant violation
of the noncontextual budget on the even fabric, i.e. an estimate consistent with $0$ to the same few-
percent resolution.  Because the control arm has no contextual signal to integrate, its budget is set by
the upper-confidence target rather than by a violation; a comparable detected-event count suffices to
exclude $\mathrm{CF}\gtrsim$ a few percent.

\begin{decisionbox}
\begin{itemize}[leftmargin=1.4em]
\item \textbf{Pass (theory survives):} the $W(3)$ arm gives $\mathrm{CF}$ consistent with $1/10$ \emph{and}
the $W(2)$/$W(4)$ arm gives $\mathrm{CF}$ consistent with $0$.
\item \textbf{Kill:} contextuality appears on an even-order fabric, or the odd fabric fails to show
$1/10$ --- either breaks the parity law that the substrate identification rests on.
\item \textbf{Inconclusive:} the two arms are run on apparatus whose analyzer imbalance or loss profile
differs enough that the contrast cannot be attributed to the fabric; re-run with a shared calibration.
\end{itemize}
\end{decisionbox}

\begin{honestbox}
The control arm requires a different analyzer bank: the $W(2)$ and $W(4)$ ray families ($15$ and $85$
rays) are not the $40$ Witting tetrads, so this is a generalization of the v1 apparatus class, not a
re-run of the identical optics.  The claim is that the same \emph{kind} of single-photon
path/polarization-plus-qutrit setup realizes all three fabrics.  Building the even-order analyzers is
the engineering cost of buying the control.  As in v1, a positive odd-arm result certifies only the
finite contextual resource of the substrate, not the full physics program; the value added here is that
the even-arm null makes a faked contextuality far harder, because the apparatus must reproduce a
parity-dependent contrast it was never told about.
\end{honestbox}

\section{Repository command path}
\begin{verbatim}
python analysis/w33_audit_qscan.py           # parity law: CF=0 (even) vs 1/(q^2+1) (odd)
python analysis/w33_ovoid_construct.py       # explicit even-q ovoid = the control model
python analysis/holonet_control_arm.py       # both arms through the unmodified estimators
python analysis/w33_doily_mermin.py          # sign vs selection contextuality, separated
python analysis/w33_realization_dimension.py # q=2 has no complete-basis realization
holonet audit                                # pin the architecture ledger to the same commit
\end{verbatim}
The first line establishes the prediction, the second constructs the noncontextual model the control arm
must reproduce, and the third synchronizes the executable specification with the physical run.

\bigskip
\noindent\rule{\linewidth}{0.4pt}\\[0.25em]
{\footnotesize \textbf{Companion to:} \texttt{holonet\_demonstrator\_protocol\_v1.tex}.\quad
\textbf{Primary witnesses:} \texttt{analysis/w33\_audit\_qscan.py},
\texttt{analysis/w33\_ovoid\_construct.py}.}

\end{document}
